\documentclass{article}
\usepackage{amsmath, amsthm, amssymb}


ewtheorem{proposition}{Proposition}

ewtheorem{corollary}[proposition]{Corollary}

\begin{document}

\section*{6.5 Improper Integrals}

\begin{proposition}
If a function $f$ satisfies the hypotheses of Definition 3 and $f(x) \geq 0$ on $[a, \omega[$, then the improper integral (\ref{eq:6.71}) exists if and only if the function (\ref{eq:6.74}) is bounded on $[a, \omega[$.
\end{proposition}

\begin{proof}
Indeed, if $f(x) \geq 0$ on $[a, \omega[$, then the function (\ref{eq:6.74}) is nondecreasing on $[a, \omega[$, and therefore it has a limit as $b \to \omega$, $b \in [a, \omega[$, if and only if it is bounded.
\end{proof}

As an example of the use of this proposition, we consider the following corollary of it.

\begin{corollary}[Integral test for convergence of a series]
If the function $x \mapsto f(x)$ is defined on the interval $[1, +\infty[$, nonnegative, nonincreasing, and integrable on each closed interval $[1, b] \subset [1, +\infty[$, then the series
\[
\sum_{n=1}^{\infty} f(n) = f(1) + f(2) + \cdots
\]
and the integral
\[
\int_{1}^{+\infty} f(x) \, dx
\]
either both converge or both diverge.
\end{corollary}

\begin{proof}
It follows from the hypotheses that the inequalities
\[
f(n + 1) \leq \int_{n}^{n+1} f(x) \, dx \leq f(n)
\]
hold for any $n \in \mathbb{N}$. After summing these inequalities, we obtain
\[
\sum_{n=1}^{k} f(n + 1) \leq \int_{1}^{k+1} f(x) \, dx \leq \sum_{n=1}^{k} f(n)
\]
or
\[
s_{k+1} - f(1) < F(k + 1) \leq s_k,
\]
where $s_k = \sum_{n=1}^{k} f(n)$ and $F(b) = \int_{1}^{b} f(x) \, dx$. Since $s_k$ and $F(b)$ are nondecreasing functions of their arguments, these inequalities prove the proposition.
\end{proof}

In particular, one can say that the result of Example 1 is equivalent to the assertion that the series
\[
\sum_{n=1}^{\infty} \frac{1}{n^{\alpha}}
\]
converges only for $\alpha > 1$.

The most frequently used corollary of Proposition 3 is the following theorem.

\end{document}